\documentclass{ctexart}
\usepackage{amssymb}
\usepackage{amsmath}
\usepackage{cancel}

\begin{document}

f(x)在区间(0,1)上有定义, $\lim_{x \to 0^+}f(x)\,=\,0$,
且$\lim_{x \to {0}^+}\frac{f(x)-f(\frac{x}{3})}{x}\,=\,0$. 证明: $\lim_{x\to 0^+}\frac{f(x)}{x}=0$.


因为这里发生的是两件不同的事情：
1. 先利用代数恒等式，让中间的函数值相消。
2. 再利用三角不等式，把整个和式的绝对值放大为各项绝对值之和。
我们把这两个步骤分开看。
第一步：中间项为什么会消去？
例如取 \(n=3\)，有：
\[
\begin{aligned}
&f(x)-f\left(\frac{x}{3}\right)\\
+{}&f\left(\frac{x}{3}\right)-f\left(\frac{x}{9}\right)\\
+{}&f\left(\frac{x}{9}\right)-f\left(\frac{x}{27}\right).
\end{aligned}
\]

将同类项整理：
\[
\begin{aligned}
={}&f(x)
-\underbrace{f\left(\frac{x}{3}\right)
+f\left(\frac{x}{3}\right)}_{0}\\
&-\underbrace{f\left(\frac{x}{9}\right)
+f\left(\frac{x}{9}\right)}_{0}
-f\left(\frac{x}{27}\right).
\end{aligned}
\]

更准确地写成：
\[
\begin{aligned}
={}&f(x)
-\cancel{f\left(\frac{x}{3}\right)}
+\cancel{f\left(\frac{x}{3}\right)}\\
&-\cancel{f\left(\frac{x}{9}\right)}
+\cancel{f\left(\frac{x}{9}\right)}
-f\left(\frac{x}{27}\right)\\
={}&f(x)-f\left(\frac{x}{27}\right).
\end{aligned}
\]

这是普通的代数加减消去，与绝对值没有关系。
第二步：取绝对值后为什么不再消去？
现在对上面的恒等式取绝对值：
\[
\begin{aligned}
&\left|f(x)-f\left(\frac{x}{27}\right)\right|\\
={}&\left|\sum_{k=0}^{2}
\left[
f\left(\frac{x}{3^k}\right)
-f\left(\frac{x}{3^{k+1}}\right)
\right]\right|.
\end{aligned}
\]

根据三角不等式：
\[
\begin{aligned}
\leq{}&
\left|f(x)-f\left(\frac{x}{3}\right)\right|\\
&+\left|f\left(\frac{x}{3}\right)-f\left(\frac{x}{9}\right)\right|\\
&+\left|f\left(\frac{x}{9}\right)-f\left(\frac{x}{27}\right)\right|.
\end{aligned}
\]

注意：中间项只在取绝对值之前的代数求和中消去，取绝对值之后，每一项都变成了一个非负数，不能再相消。
最重要的理解
可以记住这个规律：
\[
\boxed{
\left|\sum_{k=1}^{n}a_k\right|
\leq \sum_{k=1}^{n}|a_k|
}
\]

- 左边：先把所有项相加，再取绝对值。
- 右边：每一项先取绝对值，再相加。
本题恰恰是利用左边的望远镜式消去得到 \(f(x)-f(x/3^n)\)，再用右边的放缩控制它。

这两个操作的先后顺序非常重要。

\end{document}